
\needspace{1\baselineskip}
\color{uniblue}{\subsection{Токовые органы пуска охлаждения трансформатора (РТПО)}}
\color{black}

\begin{enumerate}

\item
Токовые органы пуска охлаждения трансформатора контролируют токи по сторонам силового трансформатора и формируют сигнал пуска охлаждения трансформатора.

\item
В устройстве токовые органы пуска охлаждения трансформатора представлены функциональным блоком «РТПО». Функциональный блок имеет в своем составе функции контроля тока по каждой из трех сторон трансформатора 35 кВ (ТО ВН, ТО НН1, ТО НН2). 

\item
Ввод органа пуска охлаждения каждой стороны в работу осуществляется своим программным переключателем SGF1 «Ввод\_функции».

\item
РТПО могут быть выведены из работы оперативно путем активации сигнала «ОВ РТПО», дополнительно к этому существует возможность оперативного вывода функции контроля тока каждой стороны отдельно путем активации сигналов «ОВ РТПО ВН», «ОВ РТПО НН1», «ОВ РТПО НН2», что характеризуется активным состоянием сигналов «РТПО ВН: Оперативный вывод», «РТПО НН1: Оперативный вывод», «РТПО НН2: Оперативный вывод» соответственно для сторон ВН, НН1 и НН2 на выходе функции. 

\item
При превышении заданной уставки по току любой из сторон трансформатора выдается обобщенный сигнал пуска охлаждения трансформатора «РТПО\_общ: Пуск».

\item
Алгоритм работы РТПО выполнен в соответствии с рисунком \ref{rtpo:img1}. В таблице \ref{rtpo:tbl1} приведены параметры, необходимые для настройки РТПО.

\vspace{1mm}
\begin{figure}[H]
\centering
\includegraphics[width=1\textwidth,height=1\textheight,keepaspectratio]{\fbpath/11. РТПО Т 35/_latex/img/1.pdf}
\captionof{figure}{Функциональная схема логики РТПО}\label{rtpo:img1}
\end{figure}

\FloatBarrier % Все float'ы ДО этой строки будут размещены здесь.
\end{enumerate}

